\documentclass{article}
\usepackage{amsmath}
\newenvironment{linewise}{\par\begingroup\obeylines}{\endgroup\par}
\begin{document}
\section{Nested lists and protected structures}
\begin{enumerate}
\item A first item has two sentences. Its second sentence is long enough to wrap, with a preferred punctuation boundary and additional prose.
\begin{itemize}
\item A nested item has a short sentence. Another sentence contains more ordinary words, which should retain the same list layout after source formatting.
\item[Custom label] This item uses a custom label, followed by enough prose to check the surrounding indentation.
\end{itemize}
\item The final outer item refers to the table below, with a little more prose for wrapping.
\end{enumerate}

\begin{description}
\item[Description] A description entry contains text, with an additional sentence. The second sentence is deliberately longer than the soft width.
\end{description}

\begin{tabular}{ll}
  First column & Second column \\
  A, B; C: D   & Keep the table source intact. \\
\end{tabular}

\begin{align*}
  a, b &= c + d \\
  e &= f + g
\end{align*}

\begin{linewise}
This newline-sensitive environment must keep its source lines.
Another line, with punctuation; stays on this line.
\end{linewise}

%\begin{noindent}
This protected block has two sentences. Its original source wrapping stays intact.
%\end{noindent}
\end{document}
